\documentclass[12pt,a4paper]{article}

\usepackage[T1]{fontenc}
\usepackage[utf8]{inputenc}
\usepackage[brazil]{babel}
\usepackage{geometry}
\usepackage{microtype}
\usepackage{xcolor}
\usepackage{listings}
\usepackage{enumitem}
\usepackage{indentfirst}
\usepackage{tabularx}
\usepackage{booktabs}
\usepackage[hidelinks]{hyperref}

\geometry{
    top=2.5cm,
    bottom=2.5cm,
    left=3cm,
    right=3cm
}

\definecolor{codebackground}{RGB}{245,247,250}
\definecolor{codeborder}{RGB}{210,215,222}

\lstset{
    basicstyle=\ttfamily\small,
    backgroundcolor=\color{codebackground},
    frame=single,
    rulecolor=\color{codeborder},
    breaklines=true,
    breakatwhitespace=false,
    showstringspaces=false,
    columns=fullflexible,
    keepspaces=true,
    aboveskip=0.8em,
    belowskip=0.8em,
    literate={á}{{\'a}}1 {à}{{\`a}}1 {â}{{\^a}}1 {ã}{{\~a}}1 {é}{{\'e}}1 {ê}{{\^e}}1
             {í}{{\'i}}1 {ó}{{\'o}}1 {ô}{{\^o}}1 {õ}{{\~o}}1 {ú}{{\'u}}1 {ç}{{\c{c}}}1
             {Á}{{\'A}}1 {É}{{\'E}}1 {Ç}{{\c{C}}}1 {├}{{|}}1 {└}{{`}}1 {─}{{-}}1 {│}{{|}}1
}

\setlist[itemize]{noitemsep, topsep=0.4em}
\setlist[enumerate]{noitemsep, topsep=0.4em}
\setlength{\parindent}{1.25cm}
\setlength{\parskip}{0.25em}
\setlength{\emergencystretch}{3em}
\renewcommand{\arraystretch}{1.25}

\begin{document}

\hypersetup{pageanchor=false}

\begin{titlepage}
    \thispagestyle{empty}
    \centering

    \vspace*{1.5cm}

    {\large\bfseries PROCESSO SELETIVO Bry\par}

    \vspace*{\fill}

    {\Huge\bfseries RELATÓRIO DO DESAFIO TÉCNICO\par}

    \vspace{0.8cm}

    {\Large Desenvolvedor Back-End Java\par}

    \vspace*{\fill}

    {\small\bfseries CANDIDATO\par}

    \vspace{0.4cm}

    {\Large\bfseries Henrique Mateus Teodoro\par}

    \vspace{2cm}

    {\large 2026\par}
\end{titlepage}

\hypersetup{pageanchor=true}
\pagenumbering{roman}
\tableofcontents
\clearpage
\pagenumbering{arabic}

\section{Introdução}

Este relatório apresenta o desenvolvimento do desafio técnico para a vaga de Desenvolvedor Back-End Java da Bry. O desafio é composto por quatro etapas práticas: o cálculo do resumo criptográfico de um documento, a geração de uma assinatura digital, a verificação dessa assinatura e a disponibilização das operações em uma API REST. A quinta etapa é este relatório, e a sexta, opcional, é a distribuição do código por um pipeline.

Todas as etapas foram implementadas, incluindo a opcional. O projeto gera o hash SHA-512 do documento, produz uma assinatura CMS attached com SHA-512 e RSA, verifica a integridade da assinatura e a confiança do certificado na cadeia fornecida e expõe as operações de assinatura e verificação na API. A cada versão, um pipeline publica o executável, os arquivos das etapas 1 e 2 e uma imagem Docker.

Além de cumprir as etapas, o objetivo foi entregar um código robusto, organizado e fácil de executar, com testes automatizados, tratamento de erros consistente e um ambiente padronizado com Docker.

\section{Ambiente e tecnologias}

O desenvolvimento foi realizado em Linux (Ubuntu 24.04). As tecnologias utilizadas foram:

\begin{itemize}
    \item \textbf{Java 17} (17.0.20), conforme exigido pelo enunciado;
    \item \textbf{Spring Boot 3.5.16}, a versão mais recente da linha 3.x;
    \item \textbf{BouncyCastle 1.86} (\texttt{bcprov-jdk18on} e \texttt{bcpkix-jdk18on});
    \item \textbf{Maven}, com o Maven Wrapper (\texttt{./mvnw}) fixando a versão 3.9.16;
    \item \textbf{JUnit 5}, \textbf{AssertJ}, \textbf{MockMvc} e \textbf{JaCoCo} para testes e cobertura;
    \item \textbf{Docker} e \textbf{Docker Compose} para padronizar a compilação e a execução;
    \item \textbf{GitHub Actions} para integração contínua e publicação das versões.
\end{itemize}

O \texttt{pom.xml} também usa o plugin \texttt{maven-enforcer} para impedir a compilação com outra versão do Java, evitando diferenças entre ambientes.

\section{Estrutura e arquitetura do projeto}

O código está organizado em quatro pacotes, separando o núcleo criptográfico das formas de uso (API e linha de comando):

\begin{lstlisting}
src/main/java/br/com/bry/challenge/
|-- api/        Endpoints REST, resposta da verificação e tratamento de erros
|-- cli/        Execução das etapas 1, 2 e 3 pela linha de comando
|-- config/     Configurações da aplicação
`-- crypto/     Hash, PKCS12, assinatura e verificação CMS
\end{lstlisting}

As principais classes do pacote \texttt{crypto} são:

\begin{itemize}
    \item \texttt{DigestService}: cálculo de hash;
    \item \texttt{Pkcs12CredentialLoader} e \texttt{SigningCredential}: leitura da chave privada e do certificado;
    \item \texttt{CmsSigner}: geração da assinatura CMS;
    \item \texttt{CmsVerifier} e \texttt{VerificationResult}: verificação da assinatura e seu resultado;
    \item \texttt{CertificateInfo}: informações do certificado em formato legível, usadas pela linha de comando e pela API;
    \item \texttt{CertificateTrustValidator}: validação da cadeia de confiança;
    \item \texttt{CryptoException} e \texttt{CryptoErrorCode}: erros do domínio.
\end{itemize}

O mesmo núcleo é usado pela API e pela linha de comando, sem duplicação de código. A seguir estão as principais decisões de organização.

\subsection{API e linha de comando no mesmo executável}

As etapas 1, 2 e 3 são executadas por um perfil do Spring, o \texttt{cli}. Com esse perfil ativo, a aplicação não sobe o servidor web. A classe \texttt{ChallengeStepsRunner} executa as três etapas, grava os resultados e a aplicação encerra. Sem o perfil, a aplicação sobe a API. Assim, um único jar e uma única imagem Docker atendem aos dois usos, e a linha de comando reaproveita a mesma configuração e injeção de dependências da API.

\subsection{Provider criptográfico}

Uma única instância do \texttt{BouncyCastleProvider} é criada como bean do Spring e injetada nas classes que fazem as operações criptográficas, que a usam de forma explícita. Com isso, não foi necessário registrar o provider globalmente com \texttt{Security.addProvider}, o que alteraria o estado de toda a JVM.

\subsection{Tratamento de erros no domínio}

Os erros das operações criptográficas são representados por uma única exceção, \texttt{CryptoException}, acompanhada de um código do enum \texttt{CryptoErrorCode} (por exemplo, \texttt{INVALID\_PASSWORD} ou \texttt{INVALID\_SIGNATURE\_FORMAT}). Essa abordagem permite tratar cada tipo de erro de forma específica sem criar uma classe de exceção para cada caso. Na API, cada código é convertido para o status HTTP adequado.

Uma assinatura bem formada, mas inválida, não é tratada como erro. Nesse caso a verificação retorna um resultado negativo com os motivos. Exceções ficam reservadas para entradas que não podem ser processadas, como um arquivo que não é uma assinatura CMS.

\subsection{Relógio injetável}

O momento da assinatura e a data da verificação são obtidos de um \texttt{java.time.Clock} injetado. Em produção, ele usa o relógio do sistema. Nos testes, um relógio fixo torna os resultados determinísticos. Isso evita, por exemplo, que os testes passem a falhar quando o certificado do desafio expirar.

\section{Implementação das etapas}

\subsection{Etapa 1 --- Resumo criptográfico (SHA-512)}

A classe \texttt{DigestService} calcula o hash a partir de um \texttt{InputStream}, lendo o conteúdo em blocos, sem carregar o arquivo inteiro em memória. As versões que recebem um arquivo (\texttt{Path}) ou bytes em memória apenas adaptam a entrada e chamam esse mesmo método, de modo que existe um único caminho de cálculo. O algoritmo é um parâmetro, o que permite reaproveitar o serviço na verificação da assinatura com o algoritmo informado pela própria assinatura.

O resultado é gravado em \texttt{artifacts/doc.txt.sha512}, no mesmo formato do comando \texttt{sha512sum}, o que permite conferi-lo de forma independente com \texttt{sha512sum -c}. O hash obtido para o documento é apresentado abaixo, dividido em duas linhas:

\begin{lstlisting}
dc1a7de77c59a29f366a4b154b03ad7d99013e36e08beb50d976358bea7b0458
84fe72111b27cf7d6302916b2691ac7696c1637e1ab44584d8d6613825149e35
\end{lstlisting}

\subsection{Etapa 2 --- Assinatura digital}

\subsubsection{Leitura do PKCS12}

A classe \texttt{Pkcs12CredentialLoader} lê o arquivo PKCS12 com a classe \texttt{KeyStore} do Java, como pede o enunciado, e retorna um \texttt{SigningCredential} com a chave privada e a cadeia de certificados. Ela oferece dois métodos com regras explícitas:

\begin{itemize}
    \item \texttt{load(pkcs12, senha, alias)}: usa exatamente a chave do alias informado (usado pela linha de comando);
    \item \texttt{load(pkcs12, senha)}: exige que o arquivo tenha exatamente uma chave privada (usado pela API, que recebe um PKCS12 qualquer).
\end{itemize}

Se o alias não existir, ou se houver várias chaves e nenhum alias, a operação falha com uma mensagem que lista as chaves disponíveis. Também são validados a senha, o formato do arquivo e o tipo da chave, que precisa ser RSA. Uma senha incorreta é diferenciada de um arquivo corrompido, para que a mensagem de erro seja precisa.

\noindent\textcolor{red}{\textbf{Importante:}} Durante a implementação, foi identificado que o alias informado no enunciado ``\texttt{e2618a8b-20de-4dd2-b209-70912e3177f4}'' não corresponde ao alias gravado no arquivo, que está entre chaves: ``\texttt{\{e2618a8b-20de-4dd2-b209-70912e3177f4\}}''. A diferença foi confirmada com a ferramenta \texttt{keytool} do Java. A configuração da aplicação usa o alias real, presente no arquivo.

\subsubsection{Geração da assinatura}

A classe \texttt{CmsSigner} gera a assinatura com as classes do BouncyCastle:

\begin{enumerate}
    \item \texttt{JcaContentSignerBuilder} com o algoritmo \texttt{SHA512withRSA};
    \item \texttt{JcaSignerInfoGeneratorBuilder}, que monta as informações do signatário e os atributos assinados;
    \item \texttt{JcaCertStore}, que embute o certificado do signatário na assinatura;
    \item \texttt{CMSSignedDataGenerator}, com o parâmetro de encapsulamento ativado, o que torna a assinatura attached;
    \item codificação final em DER.
\end{enumerate}

O atributo \texttt{signingTime} é criado explicitamente a partir do relógio da aplicação, já que a data da assinatura é uma das informações exigidas na verificação. Os atributos \texttt{contentType}, \texttt{messageDigest} e \texttt{cmsAlgorithmProtection} são adicionados pelo BouncyCastle.

O resultado é gravado em \texttt{artifacts/doc.txt.p7s}. A assinatura foi conferida de forma independente com o OpenSSL (\texttt{openssl cms -verify}), que confirmou a assinatura e a cadeia de confiança. A estrutura ASN.1 do arquivo mostra os elementos esperados:

\begin{lstlisting}
OBJECT        :pkcs7-signedData
OBJECT        :sha512                      (algoritmo de hash)
OBJECT        :pkcs7-data
OCTET STRING  :Teste vaga back-end Java    (documento attached)
OBJECT        :contentType
OBJECT        :signingTime
UTCTIME       :261007034340Z
OBJECT        :1.2.840.113549.1.9.52       (cmsAlgorithmProtection)
OBJECT        :messageDigest
OBJECT        :sha512WithRSAEncryption     (algoritmo de assinatura)
\end{lstlisting}

O BouncyCastle 1.86 registra o algoritmo de assinatura como \texttt{sha512WithRSAEncryption}, conforme o identificador utilizado para RSA com SHA-512 em CMS.

\subsection{Etapa 3 --- Verificação da assinatura}

\subsubsection{Integridade}

A classe \texttt{CmsVerifier} recebe a assinatura em Base64 ou em DER (o Base64 tem prioridade). O conteúdo é decodificado como Base64 e, se não for um Base64 válido, é tratado como DER.

A verificação ocorre em duas fases. Na primeira, a estrutura CMS é lida e os dados necessários à verificação são extraídos: o signatário, o conteúdo, os certificados, o algoritmo e o \texttt{signingTime}. O BouncyCastle lê algumas partes da estrutura apenas quando elas são acessadas e sinaliza dados malformados com exceções não checadas. Por isso, qualquer falha dessa fase é convertida no erro \texttt{INVALID\_SIGNATURE\_FORMAT}. Também são rejeitadas assinaturas detached (sem o documento) e assinaturas com mais de um signatário. Na segunda fase, a verificação trabalha apenas com dados já decodificados.

A integridade é verificada com \texttt{SignerInformation.verify}, que valida a assinatura RSA e verifica a correspondência entre o atributo \texttt{messageDigest} e o conteúdo assinado. Quando o \texttt{signingTime} está presente, também confere se o certificado era válido na data declarada. A confiança do certificado, incluindo sua validade na data atual, é avaliada separadamente pelo \texttt{CertificateTrustValidator}, utilizando a data de referência definida pelo \texttt{Clock} da aplicação. Cada falha gera um motivo específico (documento alterado, assinatura que não corresponde à chave do certificado ou certificado fora da validade na data da assinatura).

O enunciado relaciona o hash do documento ao \texttt{encapContentInfo}. Esse campo contém o próprio documento attached, então a implementação extrai o conteúdo e calcula seu hash com o \texttt{DigestService}, usando o algoritmo indicado no \texttt{digestAlgorithm} da assinatura.

\subsubsection{Confiança do certificado}

A classe \texttt{CertificateTrustValidator} carrega os certificados de \texttt{resources/cadeia/} e valida o certificado do signatário com o algoritmo PKIX do Java (\texttt{CertPathBuilder}), usando o provider do BouncyCastle. Os certificados autoassinados da pasta são usados como âncoras de confiança (AC raiz), e os demais como ACs intermediárias. Os certificados embutidos na assinatura também podem completar o caminho, mas a confiança continua dependendo das raízes configuradas.

A cadeia é carregada na inicialização da aplicação. Se os certificados não puderem ser lidos, a aplicação não sobe, já que sem a cadeia nenhuma assinatura poderia ser validada.

Duas decisões foram tomadas nesta etapa:

\begin{itemize}
    \item \textbf{Data de referência:} a confiança é avaliada na data atual. O \texttt{signingTime} é declarado pelo próprio signatário e não é autenticado por terceiros, então não serve como referência confiável sem um carimbo do tempo. Essa decisão é compatível com o modelo de validação de assinaturas da ETSI EN 319 102-1, que trata o \textit{claimed signing time} como uma indicação temporal que, em geral, não provém de uma fonte confiável. Assim, uma assinatura feita com um certificado que já expirou deixa de ser válida, mesmo que tenha sido feita dentro da validade.
    \item \textbf{Revogação:} a consulta às listas de certificados revogados (LCR) está desabilitada, pois exigiria acesso à rede para obter as listas publicadas pela AC.
\end{itemize}

\subsubsection{Resultado e informações do signatário}

O resultado da verificação é representado pelo record \texttt{VerificationResult}. Uma assinatura é considerada válida quando está íntegra \textbf{e} o certificado é confiável. O enunciado define a verificação da integridade e a da confiança como partes da etapa 3; por isso, o resultado final considera as duas, e cada uma também é apresentada separadamente.

Na linha de comando, a etapa 3 verifica o arquivo gerado na etapa 2 e imprime as informações relevantes:

\begin{lstlisting}[basicstyle=\ttfamily\footnotesize]
Etapa 3 - Verificação da assinatura (artifacts/doc.txt.p7s)
  Integridade:            true
  Certificado confiável:  true
  Resultado:              VÁLIDA
  Data da assinatura:     2026-10-07T03:43:40Z
  Algoritmo de hash:      SHA-512
  Hash do documento:      dc1a7de77c59a29f...825149e35
  Signatário (CN):        HUB2 TESTES
  Titular:                CN=HUB2 TESTES,OU=Validado por email,...
  Emissor:                ...,CN=AC BRy Servidor Seguro v3
  Número de série:        25F
  Validade:               2021-07-21T00:00:00Z a 2029-07-21T18:22:00Z
  Caminho de certificação: 1. CN=HUB2 TESTES,...
                          2. ...,CN=AC BRy Servidor Seguro v3
                          3. ...,CN=Autoridade Certificadora Raiz BRy ... v3
\end{lstlisting}

Se a assinatura gerada não for válida, a execução termina com erro.

\subsection{Etapa 4 --- API REST}

A API é implementada pela classe \texttt{SignatureController}, que apenas converte as requisições HTTP para o domínio e reaproveita as classes das etapas anteriores. Os endpoints são exatamente os do enunciado, \texttt{POST /signature/} e \texttt{POST /verify/}, e ambos recebem \texttt{multipart/form-data}.

\subsubsection{Endpoint de assinatura}

O endpoint \texttt{/signature/} recebe três campos:

\begin{itemize}
    \item \texttt{file}: arquivo a ser assinado;
    \item \texttt{pkcs12}: arquivo PKCS12 com a chave privada e o certificado;
    \item \texttt{password}: senha do PKCS12.
\end{itemize}

A resposta contém apenas a assinatura CMS em Base64, com o tipo \texttt{text/plain}, como descrito no enunciado. A senha é recebida como um campo do formulário multipart, e nunca como parâmetro da URL, convertida para \texttt{char[]} e apagada da memória após o uso.

\subsubsection{Endpoint de verificação}

O endpoint \texttt{/verify/} recebe o arquivo assinado no campo \texttt{signature}. O enunciado pede o arquivo em Base64; como complemento, também é aceito o arquivo \texttt{.p7s} binário (DER). Em contato anterior com a equipe da Bry sobre o mesmo endpoint, foi orientado que o suporte a Base64 é obrigatório e que aceitar os dois formatos seria um complemento desejável.

A resposta contém os campos obrigatórios \texttt{status} e \texttt{infos}, além de informações adicionais em \texttt{details}:

\begin{lstlisting}
{
  "status": "VALIDO",
  "infos": {
    "signerName": "HUB2 TESTES",
    "signingTime": "2026-10-07T03:43:40Z",
    "documentHash": "dc1a7de77c59a29f...825149e35",
    "digestAlgorithm": "SHA-512"
  },
  "details": {
    "integrityValid": true,
    "certificateTrusted": true,
    "certificate": {
      "subject": "CN=HUB2 TESTES,OU=Validado por email,...",
      "issuer": "...,CN=AC BRy Servidor Seguro v3",
      "serialNumber": "25F",
      "notBefore": "2021-07-21T00:00:00Z",
      "notAfter": "2029-07-21T18:22:00Z"
    },
    "certificationPath": ["CN=HUB2 TESTES,...", "...", "..."],
    "failureReasons": []
  }
}
\end{lstlisting}

Quando a assinatura é bem formada, mas inválida (por exemplo, com o documento alterado ou com um certificado fora da cadeia confiável), a requisição foi processada corretamente e a API retorna HTTP 200 com \texttt{status} igual a \texttt{INVALIDO} e os motivos em \texttt{failureReasons}.

\subsubsection{Tratamento de erros}

A classe \texttt{ApiExceptionHandler} centraliza o tratamento de erros. Todas as respostas de erro seguem o formato Problem Details (RFC 9457), com mensagens em português e o campo \texttt{code}, que identifica o motivo de forma estável:

\begin{lstlisting}
{
  "type": "about:blank",
  "title": "Bad Request",
  "status": 400,
  "detail": "Senha do PKCS12 incorreta",
  "instance": "/signature/",
  "code": "INVALID_PASSWORD"
}
\end{lstlisting}

\begin{center}
\small
\begin{tabularx}{\textwidth}{lX}
\toprule
\textbf{HTTP} & \textbf{Situação (\texttt{code})} \\
\midrule
400 & campo ausente (\texttt{MISSING\_FIELD}), arquivo vazio (\texttt{EMPTY\_FILE}), PKCS12 inválido, senha incorreta, chave ausente ou não RSA, conteúdo que não é uma assinatura CMS e algoritmo não suportado \\
413 & arquivo acima do limite de 10 MB (\texttt{FILE\_TOO\_LARGE}) \\
415 & requisição que não é \texttt{multipart/form-data} (\texttt{UNSUPPORTED\_MEDIA\_TYPE}) \\
500 & falha ao gerar a assinatura (\texttt{SIGNATURE\_FAILED}) e erros inesperados (\texttt{INTERNAL\_ERROR}) \\
\bottomrule
\end{tabularx}
\end{center}

Em erros inesperados, os detalhes técnicos são registrados apenas no log, sem serem expostos ao cliente. Os logs da API registram as operações e seus resultados, sem incluir a senha nem os nomes dos arquivos enviados.

Os endpoints foram testados por testes de integração, com \texttt{curl} e com o Postman, recomendado no enunciado. A coleção do Postman está disponível em \texttt{docs/bry-challenge.postman\_collection.json}.

\subsection{Etapa 5 --- Relatório das atividades}

A quinta etapa é representada por este documento, que descreve a implementação, a estrutura do projeto, as instruções de execução, as decisões tomadas e as dificuldades encontradas. As instruções de uso também estão no README do repositório.

\subsection{Etapa 6 --- Distribuição do código}

A etapa opcional foi implementada com o GitHub Actions. Ao enviar uma tag de versão (por exemplo, \texttt{v1.0.0}), o pipeline executa toda a integração contínua e, se tudo passar, publica uma release com o jar e os arquivos das etapas 1 e 2, além de uma imagem Docker no GitHub Container Registry. O funcionamento é detalhado na seção~\ref{sec:distribuicao}.

\section{Compilação, execução e artefatos}

Os comandos completos estão no README. Há duas formas de executar o projeto: com Docker, sem precisar de Java instalado, ou com Java 17 e o Maven Wrapper. Os comandos devem ser executados a partir da raiz do projeto.

\subsection{Com Docker}

\begin{lstlisting}
docker compose up --build        # sobe a API na porta 8080
docker compose run --rm cli      # executa as etapas 1, 2 e 3
\end{lstlisting}

O \texttt{Dockerfile} tem duas fases. Na primeira, o projeto é compilado e testado dentro de um container com o JDK 17. Na segunda, a imagem final recebe apenas o JRE, o jar e os arquivos do desafio, e a aplicação é executada com um usuário sem privilégios de administrador. O serviço \texttt{cli} do Docker Compose grava os arquivos na pasta \texttt{artifacts/} do projeto com o usuário da máquina.

\subsection{Com Java 17}

\begin{lstlisting}
./mvnw clean verify                                    # compila e testa
./mvnw spring-boot:run                                 # sobe a API
./mvnw spring-boot:run -Dspring-boot.run.profiles=cli  # etapas 1, 2 e 3
\end{lstlisting}

\subsection{Artefatos}

\begin{itemize}
    \item \texttt{artifacts/doc.txt.sha512}: hash SHA-512 do documento (etapa 1);
    \item \texttt{artifacts/doc.txt.p7s}: assinatura CMS attached em DER (etapa 2).
\end{itemize}

Os caminhos dos arquivos, o alias e a senha do PKCS12 ficam no \texttt{application.yml} e podem ser alterados por argumentos ou variáveis de ambiente.

\section{Testes e qualidade}

O projeto possui 94 testes automatizados: 75 de unidade e 19 de integração. Os testes de integração sobem a aplicação com o Spring e são identificados pelo sufixo \texttt{IntegrationTest}.

\begin{center}
\footnotesize
\begin{tabularx}{\textwidth}{lrX}
\toprule
\textbf{Classe de teste} & \textbf{Testes} & \textbf{Principais cenários} \\
\midrule
\texttt{DigestServiceTest} & 9 & vetores do NIST, arquivo do desafio, entrada vazia, equivalência entre arquivo, bytes e stream, falhas de leitura \\
\texttt{Pkcs12CredentialLoaderTest} & 12 & alias, arquivo com várias chaves, senha incorreta, arquivo inválido ou vazio, ausência de chave, chave EC \\
\texttt{SigningCredentialTest} & 3 & cadeia imutável e cadeia vazia \\
\texttt{CmsSignerTest} & 9 & conteúdo attached, algoritmos, atributos assinados, certificado embutido, validade criptográfica, DER \\
\texttt{CmsVerifierTest} & 16 & Base64 e DER, documento e assinatura alterados, chave que não corresponde ao certificado, signatário não confiável, certificado expirado, formatos inválidos \\
\texttt{CertificateTrustValidatorTest} & 8 & cadeia da Bry, AC desconhecida, intermediária ausente ou embutida, período de validade, cadeia inválida \\
\texttt{VerificationResultTest} & 2 & nome do signatário \\
\texttt{CertificateInfoTest} & 2 & dados do certificado e e-mail legível \\
\texttt{ChallengeStepsRunnerTest} & 6 & arquivos e saída das etapas 1, 2 e 3 e falhas de configuração \\
\texttt{ApiExceptionHandlerTest} & 8 & status HTTP de cada código de erro e erro inesperado \\
\texttt{SignatureControllerIntegrationTest} & 16 & assinatura, verificação em Base64 e DER, assinaturas inválidas e erros da API \\
Demais testes de integração & 3 & inicialização da aplicação, perfil \texttt{cli} e limite de upload \\
\bottomrule
\end{tabularx}
\end{center}

Os casos de erro são construídos com certificados e arquivos PKCS12 gerados em memória durante os testes, sem adicionar arquivos binários ao repositório.

A cobertura é medida pelo JaCoCo, e o build falha se ela ficar abaixo de 90\% das linhas ou de 80\% dos branches. A cobertura atual é de 97,6\% das linhas e 91,2\% dos branches. As linhas não cobertas correspondem a tratamentos de falhas internas que não podem ser provocadas com dados reais.

\section{Integração contínua e distribuição}

\subsection{Integração contínua}

O workflow de CI é executado a cada push e possui dois jobs, executados em Linux:

\begin{enumerate}
    \item \textbf{Build e testes:} compila o projeto, executa os testes e verifica a cobertura; executa as etapas 1, 2 e 3 e confere os resultados com ferramentas independentes do projeto: \texttt{sha512sum} para o hash e \texttt{openssl cms -verify} para a assinatura e a cadeia de confiança;
    \item \textbf{Imagem Docker e smoke test:} constrói a imagem, sobe o container, aguarda a aplicação ficar saudável e chama os endpoints \texttt{/signature/} e \texttt{/verify/}, exigindo o status \texttt{VALIDO}.
\end{enumerate}

O cache das dependências do Maven e das camadas da imagem reduz o tempo das execuções seguintes, e execuções antigas da mesma branch são canceladas quando um novo commit é enviado.

\subsection{Distribuição}
\label{sec:distribuicao}

O workflow de release é disparado por tags no formato \texttt{vX.Y.Z}. Ele reaproveita o workflow de CI, sem duplicar os passos de build e teste, e só publica a versão se toda a CI passar. Em seguida:

\begin{itemize}
    \item confere se a versão da tag corresponde à versão do \texttt{pom.xml};
    \item publica uma release no GitHub com o jar testado pela CI e os arquivos das etapas 1 e 2;
    \item publica a imagem Docker no GitHub Container Registry, com a tag da versão e a tag \texttt{latest}.
\end{itemize}

Com a imagem, a API pode ser executada sem clonar o repositório:

\begin{lstlisting}
docker run -p 8080:8080 \
    ghcr.io/henrique1803/bry-java-backend-challenge:1.0.0
\end{lstlisting}

\section{Dificuldades encontradas}

A primeira dificuldade foi a divergência entre o alias do certificado informado no enunciado e o alias gravado no arquivo PKCS12. Ao usar o valor do enunciado, a chave não era encontrada. A causa foi identificada inspecionando o arquivo com o \texttt{keytool}, com o alias real estando entre chaves.

Outra dificuldade foi tornar a verificação robusta para arquivos malformados. A leitura inicial da assinatura já tratava conteúdos inválidos, mas, ao testar a verificação com assinaturas corrompidas em posições variadas, algumas falhas ainda geravam exceções não tratadas. A causa é que o BouncyCastle lê partes da estrutura, como a chave pública do certificado e os signatários, apenas quando elas são acessadas. Na API, essas falhas resultariam em erro 500. A solução foi separar a verificação em uma fase de leitura completa da estrutura, que converte qualquer falha em erro de formato, e uma fase de verificação sobre os dados já decodificados.

Também foi necessário definir a data de referência da validação do certificado. A primeira implementação usava o \texttt{signingTime}, mas, como esse atributo é declarado pelo próprio signatário, a validação passou a usar a data atual, em linha com o modelo de validação da ETSI EN 319 102-1 para assinaturas sem carimbo do tempo.

Por fim, a interpretação de alguns pontos do enunciado exigiu decisões, como o critério para o status \texttt{VALIDO}, o formato aceito pelo endpoint de verificação e o formato da resposta da assinatura. Essas decisões foram tomadas buscando seguir o enunciado e estão descritas nas seções correspondentes. A pouca experiência prévia com a biblioteca BouncyCastle também foi uma dificuldade, principalmente nas classes de geração e verificação de assinaturas CMS. Além disso, o prazo para a realização do desafio foi curto, de três dias, e coincidiu com dias de semana. Por isso, foi necessário conciliar o desenvolvimento do desafio com as tarefas do meu dia a dia.

\section{Limitações e melhorias futuras}

\begin{itemize}
    \item \textbf{Revogação:} a consulta às LCRs (ou OCSP) não foi implementada, pois exigiria acesso à rede durante a verificação e os testes.
    \item \textbf{Validade de longo prazo:} sem um carimbo do tempo, uma assinatura deixa de ser válida quando o certificado expira. Para preservar a validade, seria necessário adotar uma política com carimbo do tempo, como a AD-RT da ICP-Brasil.
    \item \textbf{Senha na API:} a senha do PKCS12 é recebida pelo Spring como \texttt{String}, que não pode ser apagada da memória. A conversão para \texttt{char[]} limita sua permanência apenas no código da aplicação.
    \item \textbf{Políticas de assinatura:} a verificação não restringe algoritmos considerados fracos (como SHA-1) nem verifica o uso de chave do certificado.
    \item \textbf{Formatos de assinatura:} assinaturas detached e com mais de um signatário não são suportadas.
    \item \textbf{Segurança da API:} a API não possui autenticação, o que seria necessário em um ambiente de produção.
\end{itemize}

O projeto não utiliza banco de dados, já que nenhuma etapa do desafio exige persistência.

\section{Repositório do projeto}

O código-fonte, os testes, os artefatos e as instruções de execução estão disponíveis no repositório:

\begin{center}
    \url{https://github.com/Henrique1803/bry-java-backend-challenge}
\end{center}

\section{Conclusão}

O desafio foi concluído com todas as etapas solicitadas, incluindo a etapa opcional de distribuição. O projeto calcula e grava o hash SHA-512 do documento, gera uma assinatura CMS attached com SHA-512 e RSA, verifica sua integridade e a confiança do certificado na cadeia fornecida e apresenta as informações do signatário.

A API REST reutiliza as mesmas classes da linha de comando e oferece os endpoints de assinatura e verificação no formato solicitado, com tratamento de erros padronizado. A implementação é validada por 94 testes automatizados, por um pipeline que também confere os resultados com ferramentas independentes e por uma imagem Docker publicada a cada versão.

\end{document}
